\documentclass[UTF8,a4paper,11pt,openany]{ctexbook}
\usepackage{../common/statlearnbook}
\SLSetBookMetadata{第 1 册}{统计学习方法讲义 v2.6.0}{数学基础与统计学习基本理论}
\begin{document}
\setcounter{page}{79}
\setcounter{chapter}{5}
\setcounter{figure}{0}
\pagestyle{headings}
\chaptermark{常用分布与统计推断}
\begin{lemma}[标准化]
若$X\sim\mathcal N(\mu,\sigma^2)$，则$Z=(X-\mu)/\sigma\sim\mathcal N(0,1)$。
\end{lemma}
\paragraph{首次阅读检查：证明：标准化。}\textbf{第一步。}先写证明目标与合法条件，标出第一个可执行数学动作。
\begin{proof}
使用变量代换$x=\mu+\sigma z$，雅可比绝对值为$\sigma$。把原密度与雅可比相乘后，$\sigma$与归一化常数中的$1/\sigma$抵消，指数项化为$-z^2/2$，因此得到标准高斯密度。
\end{proof}
% v2.6.0 figure UID: FIG-P077-01
% Image-reference reverse redraw: exact Gaussian formulas retained; comparison hierarchy and label clearance refined.
\tikzset{
  slfig-FIG-P077-01/.style={every path/.append style={line cap=round,line join=round}},
  slfig-FIG-P077-01-label/.style={
    slfig direct label,font=\fontsize{9.2pt}{11.2pt}\selectfont,
    fill=white,fill opacity=.95,text opacity=1,inner xsep=1.1pt,inner ysep=.8pt
  }
}
\pgfplotsset{slfig-FIG-P077-01-axis/.style={
  axis line style={draw=SLInk!90,line width=.60pt},
  tick style={draw=SLInk!80,line width=.55pt},
  tick label style={font=\fontsize{8.8pt}{10.5pt}\selectfont},
  label style={font=\fontsize{9.4pt}{11.2pt}\selectfont},
  clip=false
}}
\pgfkeysifdefined{/tikz/alt}{}{\tikzset{alt/.code={}}}
\begin{figure}[H]
\centering
\begin{tikzpicture}[slfig-FIG-P077-01,alt={对象—关系—结论：方差增大使高斯密度变宽、峰值降低，但曲线下面积保持为1}]
\begin{axis}[slfig-FIG-P077-01-axis,
  slfig axis,width=.86\linewidth,height=6.2cm,
  axis lines=left,xlabel={$x$},ylabel={密度},
  xmin=-4.5,xmax=4.5,ymin=-.065,ymax=.43,
  domain=-4.5:4.5,samples=321,smooth,
  xtick={-4,0,4},ytick={0,.2,.4},clip=false
]
  \path[name path=zero] (axis cs:-4.5,0)--(axis cs:4.5,0);
  \addplot[name path=wide,draw=none,no marks]
    {1/sqrt(8*pi)*exp(-x^2/8)};
  \addplot[fill=SLTeal,fill opacity=.08,draw=none]
    fill between[of=wide and zero];
  \addplot[name path=narrow,draw=none,no marks]
    {1/sqrt(2*pi)*exp(-x^2/2)};
  \addplot[fill=SLBlue,fill opacity=.09,draw=none]
    fill between[of=narrow and zero];
  \addplot[draw=SLBlue,line width=1.05pt,no marks]
    {1/sqrt(2*pi)*exp(-x^2/2)};
  \addplot[draw=SLTeal,line width=.90pt,
    dash pattern=on 3.4pt off 2.0pt,no marks]
    {1/sqrt(8*pi)*exp(-x^2/8)};
  \draw[draw=SLGray!85,line width=.55pt,dash pattern=on 2.8pt off 1.8pt]
    (axis cs:0,0)--(axis cs:0,.405);
  \node[slfig-FIG-P077-01-label,anchor=south west] at (axis cs:.95,.305)
    {$\mathcal N(0,1)$；峰值 $1/\sqrt{2\pi}$};
  \node[slfig-FIG-P077-01-label,anchor=south west] at (axis cs:2.05,.135)
    {$\mathcal N(0,2^2)$；峰值 $1/(2\sqrt{2\pi})$};
  \draw[draw=SLGray!88,line width=.52pt,decorate,
    decoration={brace,mirror,amplitude=3pt}]
    (axis cs:-3.5,-.018)--(axis cs:3.5,-.018)
    node[pos=.50,below=3pt,font=\fontsize{9.2pt}{11.0pt}\selectfont,
      fill=white,inner sep=.8pt]
    {两条曲线：面积 $=1$};
\end{axis}
\end{tikzpicture}
\caption{方差增大使高斯密度变宽、峰值降低，但曲线下面积保持为1}\label{fig:V1-C05-gaussian}
\end{figure}

比较\cref{fig:V1-C05-gaussian}中的曲线时要同时检查位置、尺度和面积：改变$\mu$只平移中心，改变$\sigma$会展宽并降低峰值；若只横向拉伸而不调整$1/(\sqrt{2\pi}\sigma)$，所得曲线便不再是概率密度。
\SLReviewSection{Gamma与狄利克雷分布}{sec:v260-fig077-context}
\end{document}

